\section{Comparative Evaluation Matrix}

The evaluation is designed around one falsifiable question: whether explicit slosh-state prediction improves liquid response beyond what can be obtained by stronger command smoothing alone.
The matrix therefore separates liquid evidence, model-side diagnostics, task completion, path tracking, command behavior, and runtime feasibility.

\begin{table}[t]
  \centering
  \scriptsize
  \setlength{\tabcolsep}{3pt}
  \caption{Metric groups for interpreting the experiments.}
  \label{tab:method_metric_groups}
  \begin{tabular}{p{0.28\linewidth} p{0.62\linewidth}}
    \toprule
    Group & Representative metrics \\
    \midrule
    External liquid response & RGB p95/max/RMS, residual oscillation after arrival \\
    Model diagnostics & $H_{\mathrm{model}}$ p95/max/RMS, model-to-vision trend consistency \\
    Task completion & success rate, timeout, arrival time \\
    Path tracking & projection error p95, maximum path deviation \\
    Command behavior & $|\omega|$ p95, acceleration RMS, command variation \\
    Runtime & solve-time median/p95, control frequency, solver failures \\
    \bottomrule
  \end{tabular}
\end{table}

The core internal ablation uses a two-factor design, shown in \cref{tab:method_internal_ablation}.
All variants share the same reference path, base constraints, path-progress \MPCC{} structure, and task definition.
Only the slosh model/cost and smoothing level are changed.

\begin{table}[t]
  \centering
  \scriptsize
  \setlength{\tabcolsep}{3pt}
  \caption{Internal ablation matrix for separating slosh prediction from smoothing.}
  \label{tab:method_internal_ablation}
  \begin{tabular}{p{0.17\linewidth} p{0.20\linewidth} p{0.18\linewidth} p{0.35\linewidth}}
    \toprule
    Variant & Slosh state/cost & Smoothing & Interpretation \\
    \midrule
    \Bzero{} & absent & weak & Base \MPCC{} reference \\
    \Bsmooth{} & absent & strong & Smooth-only explanation \\
    \Bslosh{} & present & weak & Independent slosh-state effect \\
    \Bours{} & present & strong & Combined \SPMPC{} planner \\
    \bottomrule
  \end{tabular}
\end{table}

The comparison between \Bsmooth{} and \Bzero{} estimates the effect of stronger smoothing alone.
The comparison between \Bslosh{} and \Bzero{} estimates the independent value of explicit liquid-state propagation.
The most important comparison is \Bslosh{} or \Bours{} against \Bsmooth{}, because it tests whether the slosh-aware formulation improves upon the strongest smooth-only explanation.
Finally, \Bours{} against \Bslosh{} measures the effect of combining slosh prediction with stronger smoothing.

Same-layer local-planner baselines test whether ordinary online planning already covers the task.
DWA, TEB, \MpcLocalPlanner{}, and LT-DWA are natural candidates because they output executable base commands under runtime constraints.
They should be compared with matched start and goal conditions, reference path, velocity and acceleration limits, arrival criteria, failure criteria, and logging pipeline.
Mobile-base anti-slosh references such as fixed-path velocity profiles, input shaping, and offline liquid-constrained trajectory optimization define a second comparison axis.
They are ranked quantitatively only when their assumptions and command interfaces can be reproduced under comparable platform conditions.

The optional governor and terminal-stopping mechanisms are evaluated incrementally rather than folded into the main method claim.
The governor comparison is
\begin{equation}
  \Bours{} \quad \text{vs.} \quad \Boursgov{},
\end{equation}
with matched nominal speed reference, base constraints, path, and \SPMPC{} weights.
The terminal-stopping comparison is
\begin{equation}
  \Boursnostop{} \quad \text{vs.} \quad \Boursstop{},
\end{equation}
where the primary metrics are post-arrival RGB RMS/p95, residual oscillation decay time, terminal acceleration, terminal angular-command variation, and settling time.
